%=============================================================================!
% 11.3  RIGID WATER                                                           !
%=============================================================================!
\section{Rigid water}
\label{sec:cs-water}

Holding bonds fixed is meant to lengthen the step, and how much it
lengthens it is a question of measurement. This section builds a small
box of liquid water, runs it flexible, with only its O--H bonds held, and
rigid, and compares the energy fluctuation of the three at a range of
steps by the test of \cref{sec:in-stability}.

\subsection*{A model of water}

The TIP3P model of Jorgensen and co-workers\cite{Jorgensen1983}, in the
form ASE implements it\cite{Larsen2017}, places a charge of $-0.834e$ on
each oxygen atom and $+0.417e$ on each hydrogen, gives the oxygen atoms a
Lennard-Jones energy between one another with $\sigma =
\SI{3.15061}{\angstrom}$ and $\varepsilon = \SI{6.596}{\milli\electronvolt}$
(which ASE stores as 0.1521 kilocalories per mole, a chemists' unit that
\texttt{mdlab.units} converts), and fixes the shape of the
molecule, with O--H \SI{0.9572}{\angstrom} and the H--O--H angle
\ang{104.52}, so that H--H is \SI{1.5139}{\angstrom}; ASE's own tests run
it with all three distances held by \texttt{FixBondLengths}. The charges
describe how one molecule pulls on another; the Coulomb energy of the
pairs within a molecule is left out, since the bonds stand in for it, and
would otherwise draw each hydrogen onto its own oxygen with an energy of
$k_{\mathrm e}q_{\mathrm O}q_{\mathrm H}/r = \SI{-5.2}{\electronvolt}$ at
\SI{0.9572}{\angstrom}.

The class \texttt{WaterModel} of \texttt{mdlab.water} computes the
Lennard-Jones energy between oxygen atoms, switched off smoothly between
5 and \SI{6}{\angstrom} (\cref{sec:pe-cutoffs}), and the Coulomb energy
of all the charges and their copies by the Ewald sum of
\cref{sec:md-ewald}, less the Coulomb energy of each molecule's own three
pairs. For the flexible model it adds the three springs of the spring
model of \cref{sec:ro-freedom}, $k_{\mathrm{OH}} =
\SI{48.48}{\electronvolt\per\angstrom\squared}$ along each O--H and
$k_{\mathrm{HH}} = \SI{17.66}{\electronvolt\per\angstrom\squared}$
between the hydrogens, at rest at the TIP3P shape, whose fastest
vibration, the symmetric stretch of the model, has $\omega =
\SI{0.857}{\radian\per\femto\second}$ and a period of
\SI{7.33}{\femto\second}. Its tests check the forces against differences
of the energy, for the rigid and the flexible model, and its parameters
against ASE's.

Sixty-four molecules, turned at random, are placed on a $4\times4\times4$
grid in a cube of side \SI{12.43}{\angstrom}, which gives the density of
liquid water at \SI{298.15}{\kelvin} and \SI{0.1}{\mega\pascal},
\SI{0.997}{\gram\per\centi\metre\cubed}\cite{NISTWebBook}; half the side,
\SI{6.21}{\angstrom}, exceeds the cutoff, as \cref{sec:md-image}
requires. A grid is a poor arrangement for a liquid, and the molecules
release energy as they leave it, which heats the box. The script that
prepares it therefore follows the rigid molecules for ten rounds of
\SI{1}{\pico\second} with $\dt = \SI{2}{\femto\second}$, each round
starting from new random velocities of the kind of \cref{sec:md-start},
which discard the heat the last round released,
with no drift, no part along a held bond, and a kinetic energy of
\SI{12.93}{\milli\electronvolt} for each freedom; Chapter~12 shows that
this is the share of each freedom at \SI{300}{\kelvin}. The potential
energy falls from \SI{-2.37}{\electronvolt} on the grid to
\SI{-25.70}{\electronvolt} after the tenth round, by
\SI{0.12}{\electronvolt} in that last round, a slow change whose end
Chapter~15 learns to detect; the runs below start from that state. The
flexible molecules start from the same positions, with velocities drawn
for all nine freedoms of each molecule, and are prepared for two rounds
of \SI{0.5}{\pico\second} with $\dt = \SI{0.25}{\femto\second}$.

\subsection*{The step each model allows}

\Cref{fig:cs-step}(a) runs the three models at a range of steps, flexible
water by velocity Verlet and the other two by RATTLE, and divides the
standard deviation of the total energy, the square root of the mean of its
squared differences from its mean, by that of the kinetic energy, the
yardstick of \cref{sec:in-stability}. All three ratios grow as $\dt^2$ at
short steps, with fitted slopes of 2.05 (flexible), 2.10 (O--H held) and
2.01 (rigid) over their three shortest steps. At equal steps they differ
greatly: at \SI{1}{\femto\second} the ratio is 6.7\% for flexible water,
1.1\% with the O--H bonds held and 0.23\% for rigid water.

\begin{figure}[htb]
\centering
\includegraphics{ch11_constraints/step}
\caption{The step each model of water allows: 64 molecules, flexible (by
velocity Verlet), with the O--H bonds held, and rigid (both by RATTLE);
open markers for the last two with hydrogen made three times heavier
(\cref{sec:cs-mass}). (a) The standard deviation of the total energy over
that of the kinetic energy, against the step, over \SI{1}{\pico\second}
(flexible, O--H held) or \SI{2}{\pico\second} (rigid); dotted, the 1\%
level. (b) The total energy less its starting value for rigid water at 2
and \SI{7}{\femto\second} and flexible water at
\SI{0.5}{\femto\second}.}
\label{fig:cs-step}
\end{figure}

How small the ratio must be is a judgement, depending on what is to be
measured (\cref{sec:in-stability}). Taking 1\% as the level, and reading
where each curve crosses it between measured points, the flexible model
allows \SI{0.41}{\femto\second}, the model with its O--H bonds held
\SI{0.97}{\femto\second}, and the rigid model \SI{1.98}{\femto\second}:
holding the O--H bonds buys a factor of 2.4, and holding the angle too a
factor of 4.8. With the O--H bonds held, the fastest motion left is the
bend, with a period of \SI{17.0}{\femto\second} in this model, shorter
than the \SI{20.9}{\femto\second} of the free spring model, whose bonds
stretch as it bends; with all three distances held, the molecule cannot
vibrate at all, and the fastest motions left are the molecules' turning
back and forth in the cage of their neighbours. These steps are read from
single runs of 1 or \SI{2}{\pico\second}: the rigid run at
\SI{2}{\femto\second} gives 1.02\% over \SI{2}{\pico\second} and 0.93\%
over \SI{10}{\pico\second} (\cref{sec:cs-trajectory}), and since the
ratio grows as $\dt^2$ a difference of 10\% in it moves the step at 1\% by
about 5\%, a spread that Chapter~15 turns into an error bar. \Cref{fig:cs-step}(b) shows that even at
\SI{7}{\femto\second}, where the ratio is 13\%, the energy of rigid water
stays in a band over \SI{2}{\pico\second}, its mean moving by
\SI{0.007}{\electronvolt\per\pico\second}, as a symplectic method's
should; at \SI{2}{\femto\second} the band is \SI{0.017}{\electronvolt}
wide for the whole box, \SI{0.26}{\milli\electronvolt} per molecule.

The rigid model is a different model, however long its step. Water held
rigid cannot stretch or bend, and any property that depends on those
motions, such as its spectrum at high frequencies, is lost; the
properties of liquid water that depend on how molecules move as wholes
are what a rigid model is built to give. A rigid model is therefore a
choice of model as much as a means of lengthening the step, and the
planned comparisons with learned propagators in Part~VI use the flexible
model as the reference when those vibrations are retained.

\begin{practice}
The tolerance of the constraints matters. The step test above holds the
bonds to $10^{-10}$; a tolerance of $10^{-3}$ leaves errors of
\SI{1.5e-3}{\angstrom} and an energy that drifts
(\cref{sec:cs-trajectory}). A tolerance may be stated as a relative or an
absolute error of the length or of its square, so the numbers to check in
a run are the largest error of a held distance, which \texttt{mdlab}
records as \texttt{bond\_error}, and whether a tighter tolerance changes
the energy.
\end{practice}

\begin{keyidea}
Measured on 64 molecules of TIP3P water, at an energy fluctuation of 1\%
of the kinetic energy's, holding the O--H bonds lengthens the step from
0.41 to \SI{0.97}{\femto\second}, and holding all three distances to
\SI{1.98}{\femto\second}: factors of 2.4 and 4.8, bought by removing the
fastest vibrations, and with them the motions that depended on them.
\end{keyidea}

\notebook{11}{compare the three models' energy at any step}

\begin{selfcheck}
How many steps of \SI{0.41}{\femto\second} and of
\SI{1.98}{\femto\second} cover a nanosecond?
\end{selfcheck}
